% sections/sec_usage.tex

\documentclass[../104_Subfiles.tex]{subfiles}

\begin{document}

\section{\texttt{subfiles} の使い方}

\subsection{メインファイルの書き方}

メインファイルのプリアンブルに \texttt{\textbackslash usepackage\{subfiles\}} を追加し、
本文中で \texttt{\textbackslash subfile\{パス\}} を使います。

\subsection{サブファイルの書き方}

サブファイルの先頭に以下を書きます:

\begin{verbatim}
\documentclass[../メインファイル.tex]{subfiles}
\begin{document}
  ... 内容 ...
\end{document}
\end{verbatim}

メインファイルから読み込まれるときは
\texttt{\textbackslash documentclass} 行が無視され、
メインのプリアンブルが適用されます。

\subsection{3つの方法の比較}

\begin{table}[h]
  \centering
  \begin{tabular}{llll}
    \hline
    & \texttt{\textbackslash input} & \texttt{\textbackslash include} & \texttt{subfiles} \\
    \hline
    単独コンパイル    & 不可     & 不可     & \textbf{可} \\
    プリアンブル共有  & 手動     & 手動     & \textbf{自動} \\
    改ページ          & なし     & あり     & なし \\
    \texttt{\textbackslash includeonly} 対応 & 非対応 & 対応 & 非対応 \\
    \hline
  \end{tabular}
  \caption{3つのファイル分割方法の比較}
\end{table}

\end{document}
